% fig82-raycast.tex -- Figure 15.x: ray casting. Polygon (0,0),(4,0),(5,3),(2,4),(0,3); ray to the right at y=2.
% fstyle.tex -- shared TikZ styles for the figures of this book (palette from lifeos.sty).
\tikzset{%
  cell/.style={draw=ink, minimum width=7mm, minimum height=6.5mm, inner sep=0pt, font=\footnotesize, fill=white},
  cellh/.style={cell, fill=accentsoft, draw=accent},
  cella/.style={cell, fill=amberbg, draw=amber},
  cellg/.style={cell, fill=goodbg, draw=good},
  cellx/.style={cell, fill=graybg, draw=slate, text=slate},
  vtx/.style={circle, draw=ink, minimum size=6.5mm, inner sep=0pt, font=\footnotesize, fill=white},
  vtxh/.style={vtx, draw=accent, fill=accentsoft},
  vtxa/.style={vtx, draw=amber, fill=amberbg},
  vtxg/.style={vtx, draw=good, fill=goodbg},
  vtxd/.style={vtx, draw=slate, fill=graybg, text=slate},
  lab/.style={font=\footnotesize, text=slate},
  labA/.style={font=\footnotesize, text=accent},
  labM/.style={font=\footnotesize, text=amber},
  idx/.style={font=\tiny, text=slate},
  arr/.style={-{stealth}, thick, ink},
  arrA/.style={-{stealth}, thick, accent},
  arrM/.style={-{stealth}, thick, amber},
  arrG/.style={-{stealth}, thick, good},
  edge/.style={thick, ink},
  edgeA/.style={thick, accent},
  edgeM/.style={thick, amber},
}%
% \cellrow{prefix}{x0}{y}{style}{v0,v1,...}: a row of cells named prefix0, prefix1, ... (pitch 0.75cm)
\providecommand{\cellrow}[5]{\foreach \v [count=\i from 0] in {#5} {\node[#4] (#1\i) at (#2+\i*0.75,#3) {\v};}}%
% \cellidx{prefix}{n}{start}: index labels under cells prefix0 .. prefix(n-1)
\providecommand{\cellidx}[3]{\foreach \i in {0,...,\numexpr#2-1\relax} {\node[idx, below=0pt of #1\i] {\the\numexpr\i+#3\relax};}}%
% \fitfig{<tikzpicture>}: shrink the picture to \linewidth only when it would overflow.
\ifdefined\fitbox\else\newsavebox{\fitbox}\fi%
\providecommand{\fitfig}[1]{\sbox{\fitbox}{#1}\ifdim\wd\fitbox>\linewidth\resizebox{\linewidth}{!}{\usebox{\fitbox}}\else\usebox{\fitbox}\fi}%

\begin{tikzpicture}[font=\footnotesize, x=1.1cm, y=0.95cm]
\draw[ink, very thick, fill=accent!8] (0,0)--(4,0)--(5,3)--(2,4)--(0,3)--cycle;
% inside point q
\fill[accent] (1.5,2) circle (2.6pt) node[below left, labA] {$q$};
\draw[arrA, dashed] (1.5,2)--(6.3,2);
\fill[amber] (4.667,2) circle (2.6pt) node[below right=1pt, labM] {};
% outside point p
\fill[good] (5.4,1.0) circle (2.6pt) node[below, text=good, font=\footnotesize] {$p$};
\draw[arrG, dashed] (5.4,1.0)--(6.3,1.0);
\node[font=\footnotesize, text=good, anchor=west] at (6.4,1.0) {0 crossings: outside};
\node[font=\footnotesize, text=accent, anchor=west] at (6.4,2.0) {1 crossing (odd): inside};
\node[lab, anchor=west, align=left, text width=5.6cm] at (6.4,3.4) {cast a ray from the point and count edges it crosses};
\node[labM, anchor=west, align=left, text width=5.6cm] at (6.4,-0.2) {a vertex on the ray is counted once via the half-open rule: edge $(a,b)$ counts if $a_y\le y<b_y$ or $b_y\le y<a_y$};
\end{tikzpicture}
